\begin{ledgroupsized}[r]{120mm}
\footnotesize
\pstart
\noindent\textbf{Überlieferung:}
\pend
\end{ledgroupsized}
\begin{ledgroupsized}[r]{114mm}
\footnotesize
\pstart
\parindent -6mm
\makebox[6mm][l]{\textit{L}}%
Aufzeichnung:
LH~XXXVII~5 Bl.~133.
Ein Bogen~2\textsuperscript{o};
Ein Zettel (ca.~10~x~16~cm.);		%
\href{\Scan{LH_37_05_133wz}}%
{Wasserzeichenfragment} in der Blattecke; 
Papiererhaltungsmaßnahmen.
Eine Seite auf Bl.~133~r\textsuperscript{o}, zwei Zeichnungen (davon eine gestrichen) auf Bl.~133~v\textsuperscript{o};
Textfolge: Bl.~133~v\textsuperscript{o}, Bl.~133~r\textsuperscript{o}.
\pend
\end{ledgroupsized}
\selectlanguage{latin}
\vspace{8mm}
\pstart\noindent
\normalsize
\href{\Scan{LH_37_05_133v}}%
{\lbrack133~v\textsuperscript{o}\rbrack}\ 
\pend
\nobreak
\begin{edfigure}[0pt][.5em]    	%
\includegraphics[scale=1]{%
gesamttex/edit_VIII-4/images/LH_37_05_133_d1_133v.pdf}
\edcaption{\lbrack\textit{Fig.~1}\rbrack}
\end{edfigure}
\pstart\noindent
\href{\Scan{LH_37_05_133r}}%
{\lbrack133~r\textsuperscript{o}\rbrack}
\pend
\nobreak
\begin{edfigure}[.5em][1em]    	%
\hfill
\includegraphics[scale=1]{%
gesamttex/edit_VIII-4/images/LH_37_05_133_d2_133r.pdf}
\hfill
\raisebox{3ex}{%
\includegraphics[scale=1]{%
gesamttex/edit_VIII-4/images/LH_37_05_133_d3_133r.pdf}}
\hfill
\edcaption{%
\edtext{\lbrack\textit{Fig.~2}\rbrack}{%
\lemma{\lbrack\textit{Fig.~2}\rbrack}%
\Cfootnote[nonum]{Ein auf Bl.~133~v\textsuperscript{o} befindlicher, gestrichener Entwurf zum Diagramm wird nicht wiedergegeben.}}
\hspace*{45mm}
\edtext{\lbrack\textit{Fig.~3}\rbrack}{%
\lemma{\lbrack\textit{Fig.~3}\rbrack}%
\Cfootnote[nonum]{Ein gestrichener Entwurf zum Diagramm wird nicht wiedergegeben.}}}
\end{edfigure}
\pstart\noindent
Sint puncta gravia\indexS{punctum grave} \textit{A}, \textit{B}, \textit{D} et recta \textit{AB} horizontalis secetur in \textit{C}, 
ita ut sit \textit{AC} ad \textit{BC}, ut 
\edtext{gravitas puncti \lbrack\textit{B}\rbrack, ad gravitatem puncti \lbrack\textit{A}\rbrack.}{%
\lemma{}%
\Bfootnote{%
gravitas puncti 
\bars{\textit{A} \textit{ändert Hrsg.}}, ad gravitatem puncti 
\bars{\textit{B} \textit{ändert Hrsg.}}.~\textit{L}}}
Erit \textit{C} 
\edtext{\lbrack centrum\rbrack}{%
\lemma{}%
\Bfootnote{%
punctum 
\textit{L ändert Hrsg.}}}
commune gravitatis.\indexS{centrum gravitatis commune}
Sit \textit{A} in verticali\indexS{planum verticale} \textit{AD}, \textit{B} in inclinato\indexS{planum inclinatum} \textit{DB}, 
\edtext{sitque \lbrack\textit{DB}\rbrack\ ad \lbrack\textit{AD}\rbrack}{%
\lemma{}%
\Bfootnote{%
sitque
\bars{\textit{AD} \textit{ändert Hrsg.}} ad
\bars{\textit{DB} \textit{ändert Hrsg.}}~\textit{L}}}
ut \textit{B} ad \textit{A}, seu ut \textit{AC} ad \textit{BC}. 
\pend 
\pstart 
Attollatur \textit{A} in perpendiculari 
\edtext{usque}{\lemma{}\Bfootnote{usque \textit{erg.~L}}} 
ad $\alpha$, 
\edtext{et \textit{B} deprimatur}{%
\lemma{et}%
\Bfootnote{%
\textit{(1)}~\textit{B} in
\textit{(2)}~\textit{B} deprimatur~\textit{L}}}
\edtext{tantundem}{\lemma{}\Bfootnote{tantundem \textit{erg.~L}}}
in 
\edtext{inclinata usque ad $\beta$.}{\lemma{inclinata}%
\Bfootnote{%
\textit{(1)}~$\alpha$ %
\textit{(2)}~usque ad $\beta$.~\textit{L}}} 
Jungatur recta $\alpha\beta$ ipsam \textit{AB} secans in \textit{K}, dico fore $K\beta$ ad $K\alpha$ 
\edtext{\lbrack ut\rbrack}{\lemma{}\Bfootnote{ad \textit{L ändert Hrsg.}}} 
\textit{CB} ad \textit{CA}. Est enim 
\edtext{$K\beta$ ad $K\alpha$ ut \textit{KE} ad $A\alpha$,}{\lemma{$K\beta$ ad}%
\Bfootnote{%
\textit{(1)}~\textit{KA} %
\textit{(2)}~$K\alpha$ ut %
\textit{(a)}~\textit{K} ad %
\textit{(b)}~\textit{KE} ad %
\textit{(aa)}~\textit{KE} %
\textit{(bb)}~$A\alpha$,~\textit{L}}}
sed \textit{KE} vel \textit{BF} est ad $B\beta$ vel ad $A\alpha$, 
ut \textit{AD} ad \textit{DB}, seu ut \textit{CB} ad \textit{CA}\lbrack,\rbrack\ ergo $K\beta$ est 
\edtext{\lbrack ad\rbrack}{\lemma{}\Bfootnote{ad \textit{erg.~Hrsg.}}} 
$K\alpha$ ut \textit{CB} ad \textit{CA}, quod erat demonstrandum. 
\pend 
\pstart 
Hinc sequitur centrum gravitatis commune\indexS{centrum gravitatis commune} ipsorum punctorum \textit{A} et \textit{B}, 
non attolli, sed cum ex \textit{C} transeat in \textit{K} manere in horizonte\indexS{horizon}\indexS{centrum gravitatis in horizonte manens} \textit{AB}. 
\pend 



